% Section 2.1.3 Significant Driving Requirements
% Owner: Slice 1 (issue #2). Step 0 seeds the frozen scenario-ID list and one worked
% quality-attribute scenario (AV-1) as a template; Slice 1 writes the remaining bodies.
\subsubsection{Significant Driving Requirements}
\label{sec:driving-requirements}

\paragraph{Functional requirements.}
The brief asks the system to:
\begin{enumerate}[label=FR-\arabic*,nosep,leftmargin=3.2em]
  \item record sales in real time and keep a local inventory in each store, plus a
  master inventory that covers all store and \DC stock;
  \item create a shipping order from a store to the \DC when a product's store inventory
  falls to its minimum stock quantity;
  \item provide a front-end and an exchange protocol that send alerts to the central
  inventory, to store tracking, and to a variable number of other clients;
  \item support product families nested to any depth, where a product may belong to more
  than one family and the family-value report counts it once per family it is in;
  \item give managers lowest-stock, overstocked-item and product-family total-value
  reports;
  \item connect to the external POS system and, as a future capability, to an external
  analytics client, with the exchanged information defined.
\end{enumerate}

\paragraph{Quality-attribute requirements.}
The quality-attribute requirements that drive the architecture are recorded as
six-part scenarios (Hawker, ``Quality Attributes, Quality Scenarios, and Design
Tactics,'' slides 8--9). Each scenario has a frozen identifier. The brief states
availability, performance, scalability, security and usability. We found the hidden
requirements for modifiability, integrability and testability in its variable clients,
nested families and external systems. We dropped portability on the instructor's
direction. Where the brief gives no number, the response measure is our own target, which
we will confirm with the stakeholders. The scenarios use problem-level terms: the store
system, the central system, the store-to-\DC link and the managers. Section~\ref{sec:views}
names the parts of our solution.

\begingroup
\footnotesize
% \scenariotable is not printed, so each scenario declares its scn: label here.
\newcommand{\scenario}[1]{\phantomsection\label{scn:#1}\qascenario{#1}}
\scenario{AV-1}
  {A fault in the WAN link between a store and the \DC, which is part of the system's own deployment (Section~\ref{sec:allocation-view}).}
  {An omission fault: the store loses its WAN link to the \DC.}
  {The store-to-\DC link and the store system.}
  {Normal operation.}
  {The store system keeps serving reads and updates against the store inventory, logs the fault, and keeps each outbound event for delivery when the link returns.}
  {The store inventory stays readable and updatable at least 99\% of the time; this single outage lasts no more than 30 minutes; no buffered event is lost when the link returns.}

\scenario{AV-2}
  {The store-to-\DC link or the central system (transient fault).}
  {An event, or the confirmation of its receipt, is lost between a store and the central system.}
  {The store system, the store-to-\DC link and the central system.}
  {Normal operation, 3 to 10 stores at full rate.}
  {The store system keeps the event and sends it again with its original identifier until the central system confirms it. The central system discards duplicates by event identifier.}
  {100\% of store events reach the \DC; 0 events lost; 0 events applied twice.}

\scenario{AV-3}
  {An internal software fault at the \DC.}
  {The process that keeps the master record crashes.}
  {The master record and its database.}
  {Normal operation; stores keep sending events.}
  {The central system keeps accepting and storing events. The process restarts and applies the events that it missed.}
  {0 events lost; the master record is current again within 30 minutes; store operation stops for 0 seconds.}

\scenario{PERF-1}
  {The POS system of one store.}
  {Cashiers' scans make sales events at 100 events per second for a sustained period.}
  {The store system.}
  {Normal operation, peak trading.}
  {The store system applies each event to the store inventory and keeps it for delivery to the \DC.}
  {100 events per second sustained; each event is in the store inventory within 1 second.}

\scenario{PERF-2}
  {The store systems of ten stores.}
  {All stores send events at full rate: 1{,}000 events per second at the \DC.}
  {The central system and its consumers.}
  {Maximum scale: ten stores and one \DC.}
  {The central system stores each event, confirms receipt to the store, and delivers the event to every central consumer. Each consumer confirms that it received the event.}
  {1{,}000 events per second sustained; each store gets its confirmation within 1 second; an event counts as delivered when every consumer has confirmed it, and each event is delivered within 5 seconds; no backlog at the \DC grows.}

\scenario{PERF-3}
  {A \DC manager.}
  {Asks for the total value of a product family with nested sub-families.}
  {The reporting function and the master record.}
  {1{,}000 events per second arrive at the \DC.}
  {The system goes through the family tree and returns the report.}
  {The report returns within 10 seconds; the central system still sustains 1{,}000 events per second.}

\scenario{PERF-4}
  {The end of a store-to-\DC outage.}
  {The link returns, and the store has a 30-minute backlog of buffered events.}
  {The store system, the store-to-\DC link and the central system.}
  {Immediately after a 30-minute store-to-\DC outage.}
  {The store system delivers the buffered backlog to the \DC and continues to record live sales.}
  {The buffered backlog is delivered to the \DC within 10 minutes while live sales continue to be recorded at the normal rate; 0 events lost.}

\scenario{SEC-1}
  {An unauthenticated person or process, inside or outside the network.}
  {Tries to read or change a store or master inventory database.}
  {The store and master inventory databases and their access paths.}
  {Normal operation.}
  {The system rejects all access without the credentials of an authorized manager or system process, and logs the attempt.}
  {100\% of attempts rejected and logged; 0 records read or changed.}

\scenario{SEC-2}
  {An attacker on the WAN.}
  {Captures or changes messages between a store and the \DC.}
  {The store-to-\DC link.}
  {Normal operation.}
  {The system encrypts each message and checks its integrity. The central system rejects a changed message and does not confirm it, so the store system sends it again.}
  {0 events readable by the attacker; 100\% of changed messages rejected; 0 changed events applied.}

\scenario{SEC-3}
  {An authenticated store manager.}
  {Tries to read or change the inventory of a different store.}
  {The access control of the store system and the central system.}
  {Normal operation.}
  {The system denies the request and logs it.}
  {100\% of out-of-scope requests denied and logged; 0 records changed.}

\scenario{SCAL-1}
  {Chain management.}
  {The chain adds stores, from three to ten.}
  {The store systems and the central system.}
  {Normal operation; stores are added one at a time.}
  {We deploy one configured store system for each new store. It sends events to the central system like the existing stores.}
  {0 code changes to the central system or its clients; \scnref{PERF-2} still holds at ten stores.}

\scenario{SCAL-2}
  {A client developer.}
  {A new alerting client connects to the central system at runtime.}
  {The central system.}
  {Normal operation at full load.}
  {The central system registers the client and sends it each new event.}
  {0 central system restarts; 0 events lost to the existing clients; the new client receives the next event.}

\scenario{USE-1}
  {A store manager who is new to the system.}
  {Must confirm a shipping order that the store system raised.}
  {The store system and its manager screen.}
  {Normal operation.}
  {The store system shows the order that it created when a product fell to its minimum stock quantity. The manager checks the order and confirms it.}
  {After at most 1 hour of training, 9 of 10 managers confirm an order without help in less than 5 minutes.}

\scenario{USE-2}
  {A \DC manager who is new to the system.}
  {Reads the lowest-stock, overstocked-item and family-value reports.}
  {The reporting function.}
  {Normal operation.}
  {Each report shows its type, scope and units with the data.}
  {After at most 1 hour of training, managers answer 9 of 10 questions about the reports correctly.}

\scenario{MOD-1}
  {A developer.}
  {Adds a new client type that receives inventory alerts.}
  {The new client, the exchange protocol and the central system.}
  {Design time.}
  {The developer writes the client against the exchange protocol and connects it to the central system.}
  {0 changed modules outside the new client; done in at most 5 person-days.}

\scenario{MOD-2}
  {A \DC manager or a developer.}
  {The product-family hierarchy gains one more nesting level.}
  {The product-family data and the reporting function.}
  {Runtime, as a data change.}
  {The system accepts the new level as data; the reports include it.}
  {0 code changes; 100\% of family-value reports correct for the new level.}

\scenario{INT-1}
  {An analytics client developer.}
  {Connects an external analytics client over the exchange protocol.}
  {The exchange protocol and the central system.}
  {Development time.}
  {The developer writes the client against the exchange protocol and connects it, with no change to existing parts.}
  {0 changes to existing parts; done in at most 10 person-days.}

\scenario{INT-2}
  {Client developers.}
  {Add a variable number of alerting clients to the central system.}
  {The central system and its delivery to clients.}
  {Normal operation at full load.}
  {Each client connects over the exchange protocol; the central system sends events to all of them.}
  {With up to 10 clients, \scnref{PERF-2} still holds; 0 code changes to the central system.}

\scenario{TEST-1}
  {A tester.}
  {Drops 10\% of the messages and confirmations between stores and the \DC during a load test.}
  {The store system, the store-to-\DC link and the central system.}
  {Test environment, 100 events per second per store.}
  {The test harness controls the loss and compares the sent events with the applied events by event identifier.}
  {The test finds 100\% of lost or duplicate events; one test run takes at most 1 hour.}

\scenario{TEST-2}
  {A tester.}
  {Stops a store system or cuts the WAN link to measure detection and failover.}
  {The store system and its failure detection.}
  {Test environment.}
  {The harness records the time of the failure, the detection and the restored access.}
  {Recovery time measured to within 1 second; the test shows an outage of at most 30 minutes.}
\endgroup

\paragraph{Utility tree.}
Table~\ref{tab:utility-tree} is our draft utility tree, in the tabular form of
Bass, Clements, and Kazman, Table 19.1 (284--286; Hawker, ``Utility Trees,'' slides 2,
3 and 6). Each scenario has a rating of (business value, technical risk) on the
book's scale: value H is must-have, M is important but not fatal if omitted, and L is
nice to have; risk H keeps the architect awake, M is concerning, and L is confident.

\begin{table}[htbp]
\centering
\small
\begin{tabularx}{\linewidth}{|l|l|X|c|}
\hline
\textbf{Quality attribute} & \textbf{Refinement} & \textbf{Scenario} & \textbf{(V, R)}\\\hline
Availability & Local operation during an outage & \scnref{AV-1} WAN loss at a store & (H, M)\\\cline{2-4}
 & Guaranteed delivery & \textbf{\scnref{AV-2} Lost store event} & \textbf{(H, H)}\\\cline{2-4}
 & \DC-side recovery & \scnref{AV-3} Master inventory crash & (H, M)\\\hline
Performance & Per-store throughput & \scnref{PERF-1} One store at 100 events per second & (H, M)\\\cline{2-4}
 & Aggregate throughput at the \DC & \textbf{\scnref{PERF-2} Ten stores at the \DC} & \textbf{(H, H)}\\\cline{2-4}
 & Report latency under load & \scnref{PERF-3} Nested-family report under load & (M, M)\\\cline{2-4}
 & Recovery after an outage & \scnref{PERF-4} Backlog catch-up after an outage & (H, M)\\\hline
Security & Data access control & \scnref{SEC-1} Unauthorized database access & (H, L)\\\cline{2-4}
 & Channel confidentiality and integrity & \scnref{SEC-2} Intercepted store--\DC traffic & (H, M)\\\cline{2-4}
 & Authorization scope & \scnref{SEC-3} Manager out of scope & (M, L)\\\hline
Scalability & Store growth & \scnref{SCAL-1} Three to ten stores & (H, M)\\\cline{2-4}
 & Client growth & \scnref{SCAL-2} Client connects at runtime & (M, L)\\\hline
Usability & Training effort for orders & \scnref{USE-1} Confirming a shipping order & (H, L)\\\cline{2-4}
 & Training effort for reports & \scnref{USE-2} Three report types & (H, L)\\\hline
Modifiability & New client type & \scnref{MOD-1} New client type & (M, L)\\\cline{2-4}
 & Product-family structure & \scnref{MOD-2} Extra nesting level & (M, M)\\\hline
Integrability & External clients & \scnref{INT-1} External analytics client & (M, M)\\\cline{2-4}
 & Alerting clients & \scnref{INT-2} Variable alerting clients & (M, L)\\\hline
Testability & Delivery under loss & \scnref{TEST-1} Delivery under message loss & (M, M)\\\cline{2-4}
 & Failover timing & \scnref{TEST-2} Failover timing & (M, M)\\\hline
\end{tabularx}
\caption{Draft utility tree: (business value, technical risk).}
\label{tab:utility-tree}
\end{table}

The (H, H) scenarios are \scnref{AV-2} and \scnref{PERF-2}. Both are on the
store-to-\DC delivery path, so that path is the vertical slice of this draft. This draft
also addresses the (H, M) scenarios (\scnref{AV-1}, \scnref{AV-3}, \scnref{PERF-1},
\scnref{PERF-4}, \scnref{SEC-2}, \scnref{SCAL-1}) in Sections~\ref{sec:solution-bg}
and~\ref{sec:views}. The tree focuses our design effort and is an input to the later ATAM.
